\section{Introduction}
\label{sec:intro}

Every time we fine-tune an instruction-tuned language model on a new task, it loses some of what it learned before.
We call this loss the \emph{forgetting tax}.
Practitioners pay it routinely: models are updated in sequence for new tasks, new customers, and new data, and each update silently erodes earlier skills~\cite{mccloskey1989catastrophic,luo2023empirical}.

To plan around the tax, we need to know how it is charged.
If erosion has a characteristic shape over training time, we can schedule evaluation and replay around it.
If some kinds of skill erode faster than others, we can protect those first.
If skills learned early are more fragile, we can choose the task order.
A single ``average forgetting'' number answers none of these questions.

Existing continual instruction-tuning benchmarks mostly report that single number.
They typically give an end-of-sequence average and a backward-transfer score under one or two fixed task orders~\cite{wang2023trace,zhang2023citb,scialom2022fine,luo2023empirical}.
This leaves three things unmeasured.
First, the \emph{shape} of the forgetting trajectory between task boundaries.
Second, whether \emph{task type} matters independently of where the task sat in the order: with a fixed order, ``which task type forgets more'' and ``which task came first'' are the same comparison.
Third, whether cross-task forgetting ``rates'' survive a change of \emph{normalization}, given that different tasks are scored with different metrics.

We address these gaps by testing a concrete, three-part hypothesis.
{\bf H1 (shape):} retained performance on an old task decays as a power law in the amount of subsequent training, not uniformly.
{\bf H2 (domain):} reasoning tasks (math, logic) lose performance 3--5$\times$ faster than pattern-matching tasks (translation, summarization).
{\bf H3 (position):} early-learned tasks are forgotten disproportionately, beyond what their longer exposure to interference explains.
We fine-tune \llama and \qwen with a single LoRA adapter on six tasks in sequence.
We use a Williams balanced Latin-square design, in which every task occupies every position once and follows every other task once, and we evaluate held-out loss on all tasks every 25 optimizer steps.
We fixed the analysis plan before any training run and label every later addition as exploratory.

None of the three parts holds as stated.
Forgetting is clearly not uniform, but individual trajectories do not look like power laws: each task switch produces a hit that is largely complete within 20--40 optimizer steps and then stays flat for the rest of the 250-step stage (see \figref{fig:onset}).
On \llama, a median 55\% of a stage's loss rise has already happened at the first evaluation, 25 steps in, where uniform degradation predicts 10\%.
A power law is the best of four candidates for the \emph{averaged} \llama curves, but it clearly beats an equal-parameter saturating exponential in 0 of 36 individual trajectories.
The reasoning/pattern forgetting ratio is 1.0--1.5 on \llama depending on the normalization; on \qwen the metric-based ratios are 2.2--2.6 with wide intervals.
With lag held fixed, serial position has no measurable effect on \llama.
What does predict forgetting is the specific pair: a handful of (old task, new task) combinations carry most of the damage.

In summary, we make the following contributions:
\begin{itemize}[leftmargin=*,itemsep=0pt,topsep=0pt]
    \item We measure dense forgetting trajectories for two small instruction-tuned models under a balanced order design that separates task identity from learning position (18 sequential runs, plus zero-shot and joint-training baselines).
    \item We show that forgetting is front-loaded at each task switch and saturates within tens of steps, and that a power law describes averaged curves better than it describes the trajectories being averaged.
    \item We show that the reasoning/pattern forgetting ratio ranges from 0.88 to 1.47 on \llama and from 1.2 to 5.9 on \qwen depending only on the normalization, and that no measure supports a robust 3--5$\times$ gap.
    \item We show that the end-of-sequence disadvantage of early-learned tasks on \llama is explained by their longer exposure to interference, and that pair structure and loss of output format account for much of the observed erosion.
\end{itemize}

We review related work in \secref{sec:related}, describe the design in \secref{sec:method}, report results in \secref{sec:results}, and discuss implications and limitations in \secref{sec:discussion}.
